\begin{solution}
In the parent rest frame the two spatial momenta are opposite.  From
\((P-p_1)^2=-m_2^2\),
\[
E_1=\frac{M^2+m_1^2-m_2^2}{2M},\qquad
E_2=\frac{M^2+m_2^2-m_1^2}{2M}.
\]
Consequently
\[
|\bm p_1|=|\bm p_2|
=\frac{\sqrt{\lambda(M^2,m_1^2,m_2^2)}}{2M}.
\]
Reality gives the threshold \(M\geq m_1+m_2\).
\end{solution}

\begin{solution}
Momentum conservation is \(p+k=p'+k'\).  Using
\(p^2=p'^2=-m^2\), \(k^2=k'^2=0\), and \(p=(m,\bm0)\), squaring
\(p+k-k'=p'\) gives
\[
m(\omega-\omega')=\omega\omega'(1-\cos\theta).
\]
Thus
\[
\frac1{\omega'}-\frac1\omega=\frac{1-\cos\theta}{m},
\qquad
\omega'=\frac{\omega}{1+(\omega/m)(1-\cos\theta)}.
\]
\end{solution}

\begin{solution}
The four-dimensional measure, delta function, and positive-energy condition
are invariant.  Integrating over \(p^0\) uses
\[
\delta[-(p^0)^2+\bm p^2+m^2]\Theta(p^0)
=\frac{\delta(p^0-E_{\bm p})}{2E_{\bm p}},
\]
leaving \(\dd^3p/(2E_{\bm p})\).  Hence
\(\dd N=f(x,p)\,\dd^3x\,\dd^3p\) is expressed covariantly using the
mass-shell measure and a hypersurface element; \(f\) itself is a scalar.
\end{solution}

\begin{solution}
With signature \((-+++)\), the two independent invariants are proportional to
\[
\mathcal I_1=\bm B^2-\bm E^2,\qquad
\mathcal I_2=\bm E\cdot\bm B.
\]
If \(\mathcal I_2=0\) and \(\mathcal I_1<0\), a frame exists with
\(\bm B'=0\); if \(\mathcal I_2=0\) and \(\mathcal I_1>0\), one exists with
\(\bm E'=0\).  When both invariants vanish but the field is nonzero,
\(|\bm E|=|\bm B|\) and \(\bm E\perp\bm B\); no boost removes either field.
\end{solution}

\begin{solution}
In the instantaneous rest frame, conservation for an infinitesimal emission
gives rocket recoil momentum \(M\,\dd\chi\) equal to the photon energy
\(-\dd M\).  Hence
\[
\dd\chi=-\frac{\dd M}{M},\qquad
\chi_f=\ln\frac{M_i}{M_f}.
\]
Therefore
\[
v_f=\tanh\chi_f
=\frac{M_i^2-M_f^2}{M_i^2+M_f^2}.
\]
For small mass loss, \(v_f\simeq\ln(M_i/M_f)\), the rocket equation with
exhaust speed equal to unity.
\end{solution}

\begin{solution}
The invariant energy is
\[
s=-(p_1+p_2)^2=m_1^2+m_2^2+2m_2E_1,
\qquad E_{\rm cm}=\sqrt{s}.
\]
The total laboratory momentum is
\(\sqrt{E_1^2-m_1^2}\), so
\[
v_{\rm cm}=\frac{\sqrt{E_1^2-m_1^2}}{E_1+m_2}.
\]
At high energy \(E_{\rm cm}\sim\sqrt{2m_2E_1}\); only a fraction of order
\(\sqrt{m_2/E_1}\) of the beam energy becomes center-of-momentum energy.
\end{solution}

\begin{solution}
Decompose \(\bm s=\bm s_\perp+\bm s_\parallel\) relative to \(\bm v\).
A boost gives
\[
S^0=\gamma\,\bm v\cdot\bm s,\qquad
\bm S=\bm s_\perp+\gamma\bm s_\parallel
=\bm s+\frac{\gamma-1}{v^2}(\bm v\cdot\bm s)\bm v.
\]
Using \(p^\mu=m\gamma(1,\bm v)\) verifies \(S\cdot p=0\) and
\(S^2=\bm s^2\).  Only the longitudinal component is multiplied by
\(\gamma\).
\end{solution}

\begin{solution}
For two future timelike displacements \(a^\mu,b^\mu\), the reversed
Cauchy--Schwarz inequality gives
\[
\sqrt{-(a+b)^2}\geq\sqrt{-a^2}+\sqrt{-b^2},
\]
with equality only when they are parallel.  Repeated application shows that
the straight segment has at least the proper time of any timelike broken
line with the same endpoints.  Approximating a smooth curve by such lines and
taking the limit proves the result.
\end{solution}

\begin{solution}
Photon occupation number is invariant, as are the measures
\[
\frac{\dd^3p}{E},
\qquad
\dd^3p=\nu^2\dd\nu\,\dd\Omega.
\]
It follows that
\(I_\nu/\nu^3\) is invariant.  Since \(\nu=\mathcal D\nu'\),
\[
I_\nu(\nu,\hat{\bm n})=\mathcal D^3 I'_{\nu'}(\nu',\hat{\bm n}'),
\qquad I=\mathcal D^4 I'.
\]
For relativistic motion \(\mathcal D\) is largest inside an angle of order
\(\gamma^{-1}\) about the velocity, producing strong forward beaming.
\end{solution}

\begin{solution}
Write
\[
u^\mu=(\cosh\chi,\sinh\chi\,\hat{\bm n}).
\]
Restricting the Minkowski line element to the unit hyperboloid gives
\[
\dd\ell^2=\dd\chi^2+\sinh^2\chi\,\dd\Omega_2^2,
\]
the metric of hyperbolic three-space.  For two velocities,
\(-u_1\cdot u_2=\gamma_{\rm rel}=\cosh\chi_{\rm rel}\); therefore their
hyperbolic geodesic distance is \(\chi_{\rm rel}\).
\end{solution}
